\scanpage{133}

(16)\qquad (To 4)\qquad $Z$ est central dans $G$ i.e. $Z=\mathrm{Cent}(G)$, i.e l'opération de $G/G_0$ sur $\Pi$ se fait \deleted{\unsure{(sur)}} via un hom canonique
\[
G/G_0\longrightarrow Z\qquad [\,\text{resp. }Z/\{1,\unsure{\underline{\circ}|\underline{-}|\varepsilon}\}\ \text{dans le cas bipyramidal}\,]
\]

% slanted
\margin{\uline{NB} Si cet axiome n'est pas vérifié on peut néanmoins introduire $G'=\mathrm{Ker}(G\longrightarrow\mathrm{Aut}(Z))=\mathrm{Centr}_G(Z)$, puis $Z'=\mathrm{Im}(G'\longrightarrow Z)\supset Z^2$, puis le diagramme canonique
\diagram{\begin{tikzcd}[ampersand replacement=\&,column sep=tiny,row sep=tiny]
1\arrow[r] \& G_0\arrow[r] \& G'\arrow[r]\arrow[d,dash] \& \underset{\subset Z}{Z'}\arrow[r] \& 1\\
\& \& Z\arrow[ur,swap,"\lambda\mapsto\lambda^2"] \& \&
\end{tikzcd}}
et définir aussi $\Pi'$ comme $\Pi'=\{A_\lambda\mid\lambda\in\Pi'\}$.}

de sorte qu'on a une mirifique suite exacte

\diagram{\begin{tikzcd}[ampersand replacement=\&,column sep=large]
(17)\qquad 1\arrow[r] \& G_0\arrow[r] \& G\arrow[r,"\chi'"] \& Z\arrow[r] \& 1\,,\\
\& \& Z=\mathrm{Cent}(G)\arrow[u,hook]\arrow[ur,swap,"\lambda\mapsto\lambda^2"] \& \&
\end{tikzcd}}
% slanted
\margin{dans le cas bipyramidal, il faut remplacer le $Z$ dans la suite exacte par $Z/\{1,\underline{\circ}\}$ \deleted{\unsure{mais cette suite est exacte}}}

qui est donc associée canoniquement à $(S,F_0)$ satisfaisant (To 1) à (To 4), ou encore à l'ens. $\Pi$ de structures triangulées sur $S$ défini par \deleted{$\unsure{P_0=|P_0}$} $(S,F_0)$.

Appelons les structures triangulées satisfaisant les axiomes (To 1) à (To 4) ``précongruentielles\inserted{ symplectiques}''. On introduit sur l'ens des \deleted{\unsure{relations|structures}}\inserted{structures} triangulées précongruentielles sur l'ens $S$ \unsure{une|la} relation d'équivalence

\[
\left\{\begin{aligned}
&(S,F_0)\ \text{équivalent à }(S,F_1),\quad\deleted{\text{s'il existe}}\deleted{\text{\unsure{ces str. sont}}}\inserted{\text{s'il existe}}\\
&\lambda\in Z_{F_0}\quad\text{tel que } F_1=\lambda(F_0),\quad\text{i.e. }\deleted{\text{s'il exist}}\inserted{\text{si}}\\
&\deleted{\lambda\in\unsure{\mathfrak S|\mathcal S}}\quad\mathrm a)\ (S,F_1)\ \text{isom. à }(S,F_0)\\
&\hphantom{\lambda\in\mathfrak S\quad}\mathrm b)\ \exists\lambda\in\mathfrak S_S,\ \text{telle que }\lambda\text{ commute à}\\
&\hphantom{\lambda\in\mathfrak S\quad\mathrm b)\ }G_0=\mathrm{Aut}(S,F_0),\quad\text{et que l'on ait}\\
&\hphantom{\lambda\in\mathfrak S\quad\mathrm b)\ }\{s,t\}\in A_{F_0}\Longleftrightarrow\{s,\lambda t\}\in A_{F_1}.
\end{aligned}\right.\tag*{(18)}
\]
% slanted
\margin{on pourrait élargir la définition au cas des structures ne satisfaisant pas (To 4), cf annotation marginale plus haut\ldots}

\deleted{Soit} Une \uline{structure triangulée précongruentielle}

% archive sheet number: 132
